\subsection{室的结构}

设 \(\mathbf{C}\) 是一个室，\(\mathfrak{M}\) 是 \(\mathbf{C}\) 的墙集，并且对于 \(\mathrm{H} \in \mathfrak{M}\) ，设 \(e_{\mathrm{H}}\) 为与 \(\mathrm{H}\) 正交、关于 \(\mathrm{H}\) 与 \(\mathbf{C}\) 位于同侧的单位向量。

命题 7. 假设群 W 是本质的且是有限的。则：

\begin{enumerate}
\def\labelenumi{(\roman{enumi})}
\item
  存在唯一的点 \(a\) 属于 \(\mathbf{E}\) 且在 \(\mathbf{W}\) 下不变。
\item
  族 \((e_{\mathsf{H}})_{\mathsf{H}\in \mathfrak{M}}\) 是 \(\mathrm{T}\) 的一组基。
\item
  室 C 是以 a 为顶点、由 T 的基 \((e_{\mathrm{H}}^{\prime})_{\mathrm{H}\in \mathfrak{M}}\) （它满足 \((e_{\mathrm{H}}|e_{\mathrm{H}'}') = \delta_{\mathrm{HH}'}\) ）所确定的开单形锥。
\item
  由第 6 号命题 4，存在一点 \(a \in E\) 在 W 下不变。设 \(t \in T\) 使得 \(t + a\) 在 W 下不变。对所有 \(w \in W\) ，
\end{enumerate}

\[
\mathrm{U} (w). t + a = w (t + a) = t + a,
\]

于是 \(\mathrm{U}(w).t = t\) ；由于 \(\mathrm{W}\) 是本质的，这蕴含 \(t = 0\) ，从而证明了 \(a\) 的唯一性。

\begin{enumerate}
\def\labelenumi{(\roman{enumi})}
\setcounter{enumi}{1}
\item
  由于 W 是本质的，按第 7 号的记号有 \(T = T_{1}\) ，且命题 5 (iv) 表明族 \((e_{\mathrm{H}})_{\mathrm{H}\in\mathfrak{M}}\) 生成向量空间 T。E 中存在在 W 下不变的点这一事实表明族 \((e_{\mathrm{H}})_{\mathrm{H}\in\mathfrak{M}}\) 是自由的（第 6 号，命题 4）。
\item
  设 \(a\) 是 \(\mathbf{E}\) 的在 \(\mathbf{W}\) 下不变的唯一点。由于 \((e_{\mathrm{H}})_{\mathrm{H}\in \mathfrak{M}}\) 是 \(\mathbf{T}\) 的一组基，且标量积是 \(\mathbf{T}\) 上的非退化双线性型，故存在唯一的基 \((e_{\mathrm{H}}^{\prime})_{\mathrm{H}\in \mathfrak{M}}\) （属于 \(\mathbf{T}\) ），使得 \((e_{\mathrm{H}}|e_{\mathrm{H}'}') = \delta_{HH'}\) 对 \(\mathrm{H},\mathrm{H}'\) 属于 \(\mathfrak{M}\) 成立。每一点 \(x\) （\(\mathbf{E}\) 中的）都可唯一地写成形式 \(x = t + a\) ，其中 \(t = \sum_{\mathrm{H}\in \mathfrak{M}}\xi_{\mathrm{H}}.e_{\mathrm{H}}'\) 且诸 \(\xi_{\mathrm{H}}\) 为实数。于是 \(x\) 属于 \(\mathbf{C}\) 当且仅当：对每个超平面 \(\mathrm{H}\in \mathfrak{M}\) ，\(x\) 位于 \(\mathbf{H}\) 的与 \(e_{\mathrm{H}}\) 相同的一侧，换言之 \((t|e_{\mathrm{H}}) = \xi_{\mathrm{H}}\) 严格为正。由此得 (iii)。
\end{enumerate}

命题 8. 假设群 W 是本质的、不可约的且无限的。则：

\begin{enumerate}
\def\labelenumi{(\roman{enumi})}
\item
  E 中没有在 W 下不变的点。
\item
  我们有 \(\operatorname{Card} \mathfrak{M} = \dim T + 1\) ，并且存在实数 \(c_{\mathrm{H}} > 0\) 使得 \(\sum_{\mathrm{H} \in \mathfrak{M}} c_{\mathrm{H}} \cdot e_{\mathrm{H}} = 0\) 。若实数 \(c_{\mathrm{H}}'\) 满足 \(\sum_{\mathrm{H} \in \mathfrak{M}} c_{\mathrm{H}}' \cdot e_{\mathrm{H}} = 0\) ，则存在实数 \(\xi\) 使得 \(c_{\mathrm{H}}' = \xi c_{\mathrm{H}}\) 对 \(\mathfrak{M}\) 中所有 H 成立。
\item
  室 C 是一个开单形。
\end{enumerate}

断言 (i) 由命题 4 得出。另一方面，由于 W 是本质的，向量 \((e_{\mathrm{H}})_{\mathrm{H}\in\mathfrak{M}}\) 生成 T。我们有 \((e_{\mathrm{H}}|e_{\mathrm{H}^{\prime}})\leqslant0\) （对于 \(H,H^{\prime}\in \mathfrak{M}\) 且 \(H\neq H^{\prime}\) ；命题 3），并且由于 W 不可约，不存在把 \(\mathfrak{M}\) 分成两个不相交子集 \(\mathfrak{M}^{\prime}\) 与 \(\mathfrak{M}^{\prime\prime}\) 的分割，使得 \(H^{\prime}\in \mathfrak{M}^{\prime}\) 与 \(H^{\prime\prime}\in \mathfrak{M}^{\prime\prime}\) 蕴含 \((e_{\mathrm{H}^{\prime}}|e_{\mathrm{H}^{\prime\prime}})=0\) 。于是可以应用第 5 号引理 5，且该引理的情形 1) 被排除；事实上，诸 \(e_{H}\) 不是线性无关的，因为 W 没有不动点。断言 (ii) 由此得证。

现在证明 (iii)。把 C 的墙编号为 \(H_{0}, H_{1}, \ldots, H_{d}\) ，并置 \(t_{m} = e_{H_{m}}\) 。由 (ii)，向量 \(t_{1}, \ldots, t_{d}\) 构成 T 的一组基，因此超平面 \(H_{1}, \ldots, H_{d}\) 有一个公共点 \(a_{0}\) ，并且存在 T 的一组基 \((t_{1}^{\prime}, \ldots, t_{d}^{\prime})\) 使得 \((t_{m}|t_{n}^{\prime}) = \delta_{mn}\) ；此外，仍由 (ii)，存在实数 \(c_{1} > 0, \ldots, c_{d} > 0\) 使得

\[
t _ {0} = - \left(c _ {1}. t _ {1} + \dots + c _ {d}. t _ {d}\right).
\]

由于向量 \(t_{0}\) 与超平面 \(H_{0}\) 正交，存在实数 c 使得 \(H_{0}\) 是 E 中形如 \(x = t + a_{0}\) 且满足 \((t|t_{0}) = -c\) 的点所成的集合。

E 的每一点都可唯一地写成形式 \(x = t + a_{0}\) ，其中 \(t = \xi_{1}.t_{1}^{\prime} + \cdots + \xi_{d}.t_{d}^{\prime}\) 且 \(\xi_{1}, \ldots, \xi_{d}\) 为实数。点 x 属于 C 当且仅当它位于 \(H_{m}\) 的与 \(t_{m}\) 相同的一侧（对于 \(0 \leqslant m \leqslant d\) ）；这可翻译为不等式 \((t|t_{1}) > 0, \ldots, (t|t_{d}) > 0\) 与 \((t|t_{0}) > -c\) ，或等价地 \(\xi_{1} > 0, \ldots, \xi_{d} > 0\) ，\(c_{1}\xi_{1} + \cdots + c_{d}\xi_{d} < c\) 。由于 C 非空，故 c \textgreater{} 0。置 \(a_{m} = a_{0} + \frac{c}{c_{m}}.t_{m}'\) （对于 \(1 \leqslant m \leqslant d\) ）；于是室 C 由 E 中形如 \(a_{0} + \sum_{m=1}^{d} \lambda_{m}.(a_{m} - a_{0})\) 且满足 \(\lambda_{1} > 0, \ldots, \lambda_{d} > 0\) 与 \(\lambda_{1} + \cdots + \lambda_{d} < 1\) 的点组成，故 C 是以 \(a_{0}, \ldots, a_{d}\) 为顶点的开单形。证毕。

注. 1) 如第 8 号中那样，把 E 与 \(E_{0} \times E_{1} \times \cdots \times E_{s}\) 等同，把 W 与 \(W_{1} \times \cdots \times W_{s}\) 等同。由命题 6，室 C 于是等同于

\[
\mathrm{E} _ {0} \times \mathrm{C} _ {1} \times \dots \times \mathrm{C} _ {s},
\]

其中 \(C_{p}\) 是 \(E_{p}\) 中相对于超平面集 \(H_{p}\) 的一个室。由命题 7 与 8，诸室 \(C_{1},\ldots,C_{s}\) 中的每一个都是开单形锥或开单形。

\begin{enumerate}
\def\labelenumi{\arabic{enumi})}
\setcounter{enumi}{1}
\tightlist
\item
  假设 W 不可约且本质。若 H 与 \(H'\) 是 C 的两面墙，则 \(m_{HH'} = +\infty\) 当且仅当 \(e_{H} = -e_{H'}\) （命题 3）。由命题 7 与 8，这只可能发生在 H 与 \(H'\) 是 C 仅有的墙且 E 的维数为 1 的情形。于是，诸 \(m_{HH'}\) 中之一为无穷的唯一情形是：E 的维数为 1 且群 W 由两个不同点所关联的反射生成（参见~§2，第 4 号）。
\end{enumerate}

在一般情形，与 W 相关联的 Coxeter 矩阵的元素都是有限的，除非 \(E_{1},\ldots,E_{s}\) 中至少有一个属于前述类型。

